\documentclass[12pt, a4paper]{article}

\usepackage{amsmath, amssymb, amsthm}
\usepackage{geometry}
\usepackage{hyperref}

\geometry{margin=2.5cm}

% Theorem environments
\newtheorem{definition}{Definition}[section]
\newtheorem{theorem}{Theorem}[section]
\newtheorem{axiom}{Axiom}[section]
\newtheorem{example}{Example}[section]

\title{\textbf{Discrete Mathematics Notes} \\ \large Set Theory \& Axiom Systems}
\author{Shahrizoda}
\date{\today}

% -----------------------------------------------
\begin{document}
\maketitle
\tableofcontents
\newpage

% ================================================
\section{Basic Set Theory}

\begin{definition}[Set]
A \textbf{set} is a well-defined collection of distinct objects, called \textbf{elements}.
We write $x \in A$ to mean "$x$ is an element of $A$", and $x \notin A$ otherwise.
\end{definition}

\begin{example}
$A = \{1, 2, 3\}$, \quad $\mathbb{N} = \{0, 1, 2, 3, \ldots\}$, \quad $\emptyset = \{\}$
\end{example}

% -----------------------------------------------
\subsection{Set Relations}

\begin{definition}[Subset]
$A \subseteq B$ if every element of $A$ is also in $B$:
$$A \subseteq B \iff \forall x \, (x \in A \Rightarrow x \in B)$$
$A = B$ if and only if $A \subseteq B$ and $B \subseteq A$.
\end{definition}

\begin{definition}[Proper Subset]
$A \subsetneq B$ if $A \subseteq B$ and $A \neq B$.
\end{definition}

% -----------------------------------------------
\subsection{Set Operations}

\begin{definition}[Union]
$$A \cup B = \{ x \mid x \in A \text{ or } x \in B \}$$
\end{definition}

\begin{definition}[Intersection]
$$A \cap B = \{ x \mid x \in A \text{ and } x \in B \}$$
\end{definition}

\begin{definition}[Difference]
$$A \setminus B = \{ x \mid x \in A \text{ and } x \notin B \}$$
\end{definition}

\begin{definition}[Complement]
Given a universal set $U$:
$$A^c = U \setminus A = \{ x \in U \mid x \notin A \}$$
\end{definition}

\begin{definition}[Power Set]
The power set of $A$ is the set of all subsets of $A$:
$$\mathcal{P}(A) = \{ S \mid S \subseteq A \}$$
If $|A| = n$, then $|\mathcal{P}(A)| = 2^n$.
\end{definition}

\begin{definition}[Cartesian Product]
$$A \times B = \{ (a, b) \mid a \in A, \, b \in B \}$$
\end{definition}

% -----------------------------------------------
\subsection{Key Identities}

\begin{theorem}[De Morgan's Laws]
For any sets $A$ and $B$:
$$(A \cup B)^c = A^c \cap B^c$$
$$(A \cap B)^c = A^c \cup B^c$$
\end{theorem}

\begin{theorem}[Distributive Laws]
$$A \cap (B \cup C) = (A \cap B) \cup (A \cap C)$$
$$A \cup (B \cap C) = (A \cup B) \cap (A \cup C)$$
\end{theorem}

% ================================================
\section{Axiom Systems — ZFC}

ZFC stands for \textbf{Zermelo--Fraenkel set theory with the Axiom of Choice}.
It is the standard foundation for modern mathematics.

\begin{axiom}[Extensionality]
Two sets are equal if and only if they have exactly the same elements:
$$\forall A \, \forall B \, (\forall x \, (x \in A \iff x \in B) \Rightarrow A = B)$$
\end{axiom}

\begin{axiom}[Empty Set]
There exists a set with no elements:
$$\exists A \, \forall x \, (x \notin A)$$
This set is unique by Extensionality; we call it $\emptyset$.
\end{axiom}

\begin{axiom}[Pairing]
For any two sets $a$ and $b$, there exists a set $\{a, b\}$:
$$\forall a \, \forall b \, \exists A \, \forall x \, (x \in A \iff x = a \lor x = b)$$
\end{axiom}

\begin{axiom}[Union]
For any set $A$, there exists a set containing all elements of its elements:
$$\forall A \, \exists U \, \forall x \, (x \in U \iff \exists B \, (B \in A \land x \in B))$$
This gives $U = \bigcup A$.
\end{axiom}

\begin{axiom}[Power Set]
For any set $A$, there exists a set of all its subsets:
$$\forall A \, \exists P \, \forall B \, (B \in P \iff B \subseteq A)$$
This gives $P = \mathcal{P}(A)$.
\end{axiom}

\begin{axiom}[Infinity]
There exists an infinite set:
$$\exists A \, (\emptyset \in A \land \forall x \in A \, (x \cup \{x\} \in A))$$
This constructs the natural numbers: $0 = \emptyset$, $1 = \{0\}$, $2 = \{0,1\}$, \ldots
\end{axiom}

\begin{axiom}[Schema of Separation]
For any set $A$ and property $\varphi(x)$, there exists a subset of $A$ satisfying $\varphi$:
$$\forall A \, \exists B \, \forall x \, (x \in B \iff x \in A \land \varphi(x))$$
Written: $B = \{ x \in A \mid \varphi(x) \}$.
\end{axiom}

\begin{axiom}[Replacement]
The image of a set under any definable function is a set.
\end{axiom}

\begin{axiom}[Regularity / Foundation]
Every non-empty set $A$ contains an element disjoint from $A$:
$$\forall A \, (A \neq \emptyset \Rightarrow \exists x \in A \, (x \cap A = \emptyset))$$
This prevents sets like $A = \{A\}$ (a set containing itself).
\end{axiom}

\begin{axiom}[Choice]
For any set of non-empty sets, there exists a choice function:
$$\forall A \, (\emptyset \notin A \Rightarrow \exists f : A \to \bigcup A \, \forall B \in A \, (f(B) \in B))$$
\end{axiom}

% ================================================
\section{My Notes \& Questions}

% Write your own thoughts, questions, and observations here:

\begin{itemize}
    \item 
    \item 
    \item 
\end{itemize}

\end{document}
